\section{Exact eigenfunctions and a fixed metric}\label{sec:global}

We combine the packets with the round beam and choose a real sample
with a nonvanishing first jet and a large sign certificate.
Quantitative control of the first jet permits exact correction of the
coefficients; the certificate preserves the nodal lower bound under later
perturbations. We then make the eigenvalue simple in the remaining
round collar and use spectral persistence to pass to one smooth pair.

We use the common neighborhood, regions $V\Subset\Omega$, and exponent
$\gamma>0$ supplied by Proposition~\ref{prop:envelope}. Shrink the
neighborhood of positive pairs, if necessary, so that its associated
metrics belong to the neighborhood required there. Throughout a single
step the input pair is fixed. Constants may depend on this pair, on a
fixed positive collar width, and on fixed finite approximation orders.
Their frequency exponents, when specified below, do not have these
dependencies.

\subsection{A weighted Gaussian construction}

Fix $l_-,l_+$ as in Proposition~\ref{prop:step}, put $w=l_+-l_->0$,
and choose
\[
 s_i=l_-+iw/5\quad(1\le i\le4).
\]
For the allowed frequency $k=\sqrt{N(N+3)}$, define
\[
 b_N(t)=\frac{N}{2k}\log(1-t),\qquad t<1.
\]
Apply Proposition~\ref{prop:envelope} and choose the additive constant
in $h$ so that $h(s_2)=b_N(s_2)$. Here the notation $h(s_2)$ is
unambiguous because $h=-a_*t+\text{constant}$ on $t\le1/2$.
The constants added in this way are bounded, since
$s_2\in[1/4,1/2]$ and $N/k\le1$. For the fixed sufficiently small
seed slope $a_*$ and all sufficiently large $N$,
\[
 \frac{\dd}{\dd t}(h-b_N)
 =-a_*+\frac{N}{2k(1-t)}\ge c_0>0
 \qquad(0\le t\le1/2),
\]
where $c_0$ can be fixed independently of the collar. Consequently,
\begin{equation}\label{eq:gluing-gaps}
 h\le b_N-cw\quad(t\le s_1),\qquad
 h\ge b_N+cw\quad(s_3\le t\le s_4)
\end{equation}
for a fixed $c>0$.

Choose a smooth cutoff $0\le\beta\le1$ equal to one for $t\le s_3$
and zero for $t\ge s_4$. Set
\begin{equation}\label{eq:gluing-weight}
 W=e^{kh}+\beta(t)e^{kb_N(t)}.
\end{equation}
The second summand is extended by zero outside $t<s_4$; this is
smooth. The input pair is round on the support of $\beta$. Thus the
only residual of $\beta u_*$ is produced by derivatives of $\beta$,
where the second inequality in \eqref{eq:gluing-gaps} applies. For
every fixed nonnegative integer $q$ and every fixed $M>0$,
\begin{equation}\label{eq:beam-weighted-bounds}
 \begin{split}
 |D^q(\beta u_*)|&\le C_q k^qW,\\
 |D^q(L+k^2)(\beta u_*)|&\le C_{q,M}k^{-M}W,\\
 |\dd W|/W&\le Ck^2 .
 \end{split}
\end{equation}
All derivatives here and below are in the fixed finite coordinate
systems. To justify the first bound, work on $t\le s_4<1/2$ and use
$u_*=(1-t)^{N/2}\cos(N\vartheta)$ in finitely many angular charts.
Each derivative costs at most a fixed constant times $k$. In the
cutoff region the exponential gap absorbs the cutoff derivatives.
The second bound follows in the same way from the commutator with
$\beta$. Finally, $|\dd h|\le Ck^{2\delta}$ and the same gap control
the derivative of \eqref{eq:gluing-weight}, giving its last bound.
For a fixed collar, all derivatives of $\beta$ are fixed constants;
more quantitatively they may be bounded by $C_qw^{-q}$.

Choose a net in $\{t\ge s_1\}$ of mesh at most a fixed small multiple
of $k^{-10}$, with at most $Ck^{40}$ centers. At every center take the
fixed finite family of Proposition~\ref{prop:packets}. Denote the
resulting collection, counting packets with their center and family
index, by $\mathcal W_k$. Its cardinality is at most $Ck^{40}$.
For each $v\in\mathcal W_k$ take two independent standard real
Gaussian variables $g_v^{(1)},g_v^{(2)}$, all independent, and set
\begin{equation}\label{eq:gaussian-sum}
 u=\beta u_*+
 \sum_{v\in\mathcal W_k}
 \bigl(g_v^{(1)}\operatorname{Re}v+
       g_v^{(2)}\operatorname{Im}v\bigr).
\end{equation}
Let $\mathcal B_k$ be the event that all these coefficients have
absolute value at most $k$. A Gaussian tail estimate gives
\begin{equation}\label{eq:coefficient-event}
 \Pp(\mathcal B_k^c)\le Ck^{40}e^{-k^2/2}\longrightarrow0.
\end{equation}

Given a finite integer $J\ge2$ and a number $M>0$, choose the packet
orders so that their residuals through order $J$ decay at least as
$k^{-M-42}e^{kh}$. This choice is made before $k$ is taken large.
The packet estimates, \eqref{eq:beam-weighted-bounds}, and the
cardinality bound imply, on $\mathcal B_k$,
\begin{equation}\label{eq:sum-weighted-bounds}
 \begin{split}
 |D^qu|&\le C_q k^{2q+41}W\qquad(0\le q\le J),\\
 |D^q\mathfrak r|&\le C k^{-M}W\qquad(0\le q\le J),\\
 \mathfrak r&=\diver_\mu(A\,\dd u)+k^2\rho u .
 \end{split}
\end{equation}
The multiplication by the fixed smooth $\rho$ is included in the
second estimate. In particular, increasing $J$ or $M$ does not change
the exponent $2q+41$ in the first estimate.

Every packet is supported within distance $Ck^{-1/3-100\delta}$ of
its center. The function $t$ has bounded differential, so for a fixed
$C$ all packets vanish on
\begin{equation}\label{eq:untouched-region}
 t<s_1-Ck^{-1/3-100\delta}.
\end{equation}
On this set $u=u_*$ and $\mathfrak r=0$.

We use the full five-dimensional jet spanning in \eqref{eq:jet-span}
to show that, with probability tending to one, $u$ and its first
derivatives never become too small simultaneously. A lower bound by
an inverse power of $k$ is sufficient:
the coefficient correction will spend a larger power of residual
accuracy to compensate for it.

\begin{lemma}[A quantitative nonvanishing first jet]
\label{lem:noncancellation}
For every fixed choice of the finite packet orders, the construction
above satisfies
\begin{equation}\label{eq:noncancellation}
 \Pp\left(\mathcal B_k\cap\left\{
 \inf_{x\in B}\frac{|u(x)|+|\dd u(x)|_{g_0}}{W(x)}
 \ge k^{-500}\right\}\right)\longrightarrow1.
\end{equation}
The exponent $500$ is independent of those orders and of the fixed
input pair and collar.
\end{lemma}

\begin{proof}
We separate two deterministic regions. If $W/e^{kh}>k^{50}$, then
$\beta=1$ for large $k$: on its transition region
\eqref{eq:gluing-gaps} instead gives $W/e^{kh}\le1+e^{-cw k}$.
In the uncut region, the cosine formula and the angular derivative
give
\[
 |u_*|+|\dd u_*|_{g_0}\ge c e^{kb_N}.
\]
On $\mathcal B_k$, the packet value and first derivative together are
bounded by $Ck^{43}e^{kh}$. The beam therefore dominates their sum
on this region and gives a lower bound $cW$ for the first jet.
The first inequality in \eqref{eq:gluing-gaps} puts all of $t<s_1$
in this region for sufficiently large $k$.

At a point in the remaining region there is a nearby net center.
The persisting real jet span in \eqref{eq:jet-span}, after converting
the normalized derivative to ordinary coordinate derivatives, gives
for the value and four derivatives of the random sum a covariance
matrix whose least eigenvalue has square root at least
\[
 c k^{-1}e^{kh}\ge c k^{-51}W\ge k^{-52}W
\]
for all sufficiently large $k$. Only one nearby packet family is
needed for this lower bound. All other families add nonnegative
covariance matrices. The deterministic beam is a mean shift and
does not affect the maximum of a Gaussian density. Hence, in any
fixed coordinate chart and for any fixed constant $C_0$, the
\emph{unconditional} probability that the normalized jet
\[
 \mathcal J(x)=W(x)^{-1}(u(x),\partial_1u(x),\ldots,\partial_4u(x))
\]
has norm at most $C_0k^{-D}$ is bounded by
\begin{equation}\label{eq:jet-small-ball}
 C k^{260-5D}.
\end{equation}
Indeed the density is at most $Ck^{260}$ and the five-dimensional
ball has volume $Ck^{-5D}$. No conditioning on $\mathcal B_k$ is used
in this estimate.

On $\mathcal B_k$, \eqref{eq:sum-weighted-bounds} through order two
and the last bound in \eqref{eq:beam-weighted-bounds} give
\[
 |D\mathcal J|\le Ck^{45}\le k^{46}
\]
for sufficiently large $k$. In each of finitely many coordinate
patches covering $B$, take a grid of mesh $k^{-D-48}$. Its total
cardinality is at most $Ck^{4D+192}$. Apply
\eqref{eq:jet-small-ball} at grid points in the Gaussian region;
on the other grid points the deterministic beam bound already holds
on $\mathcal B_k$. The union of the Gaussian bad events has
probability at most
\[
 C k^{260-5D+4D+192}=Ck^{452-D}.
\]
With $D=500$ this tends to zero. Choose the fixed grid-point
threshold $C_0k^{-D}$ large enough to account for the finitely many
coordinate norm comparisons. The interpolation error on
$\mathcal B_k$ is at most $Ck^{-D-2}$, so the grid bounds imply
$|u|+|\dd u|_{g_0}\ge k^{-D}W$ everywhere. Subtracting the bad
probabilities and \eqref{eq:coefficient-event} proves
\eqref{eq:noncancellation}. All order-dependent constants have only
increased the lower threshold on $k$.
\end{proof}

\subsection{Sign certificates with a power gain}

To obtain a fixed positive probability of opposite signs on each
wavelength ball, we need a different estimate. After normalizing
the variance of the value to one, we will show that at least one
directional derivative retains variance bounded away from zero even
when the value is fixed. The nonzero imaginary logarithmic gradients
in Lemma~\ref{lem:scaled-packets} provide this bound. Thus sign
production uses the packets' oscillation at the local wavelength,
while the preceding lemma uses their full first-jet span.

\begin{lemma}\label{lem:signs}
For every fixed choice of the packet orders in
\eqref{eq:gaussian-sum}, there exist $p_*,c_*>0$, independent of
$k$, such that for all sufficiently large allowed frequencies,
the probability that $u$ has a sign certificate in $\Omega$ of size
at least $c_*k^{1+\gamma}$ is at least $p_*$. The certificate can be
chosen from pairs of disks fixed deterministically before sampling
the Gaussian coefficients.
\end{lemma}

\begin{proof}
For $x\in V$, put
\begin{equation}\label{eq:sign-radius}
 R_x=(kQ_x)^{-1},\qquad Q_x=(1+|\nabla_G h(x)|_G^2)^{1/2}.
\end{equation}
All balls below are Euclidean coordinate balls. Their radii are at
most $k^{-1}$, so fixed enlargements lie in a common relatively
compact subset of $\Omega$ for large $k$. The envelope bounds give
$|DQ|\le Ck^{4\delta}$. If two candidate balls intersect, their
centers are at distance at most $2/k$, and their $Q$ values differ
by at most $Ck^{-1+4\delta}$. Since $Q\ge1$, the corresponding
radii are uniformly comparable.

At each fixed $k$ the candidate radii have a positive lower bound.
Choose a finite maximal disjoint family of open balls
$\{B(x,R_x):x\in\mathcal X_k\}$, with centers in $V$.
Uniformly bounded enlargements of these balls cover $V$, and $Q$ is
comparable to $Q_x$ in each enlargement. Thus
\begin{equation}\label{eq:weighted-packing}
 \int_V kQ(y)\,\dd y
 \le C\sum_{x\in\mathcal X_k}kQ_xR_x^4
 =C\sum_{x\in\mathcal X_k}R_x^3,
 \qquad
 \sum_{x\in\mathcal X_k}R_x^3\ge c k^{1+\gamma}.
\end{equation}
The last inequality uses \eqref{eq:envelope-bounds}.

Fix one ball. The beam vanishes on its radius-two enlargement for
large $k$, because $\overline V\subset\{t>1/2\}$ and $s_4<1/2$.
Define
\[
 S_x^2=\sum_{v\in\mathcal W_k}|v(x)|^2,
 \qquad U_x(Y)=u(x+R_xY)/S_x,\qquad |Y|\le2.
\]
In expressions differentiated with respect to $Y$, a packet $v$
means its composition $v(x+R_xY)$.
Lemma~\ref{lem:scaled-packets} applies to the entire packet
collection, including packets centered outside $V$ that meet this
ball. It gives $S_x\ge ce^{kh(x)}$; uniformly bounded sums of
normalized squared packet derivatives through every fixed order;
and, for every non-tail packet,
\begin{equation}\label{eq:local-imaginary-lower}
 \left|\operatorname{Im}\frac{\dd_Yv}{v}(0)\right|\ge c.
\end{equation}
Here a non-tail has its cutoff identically one on the scaled ball
and nonzero value at the center. For tails, the lemma supplies
absolute exponential bounds, even if $v(x)=0$. Polynomially many
tails consequently contribute $o(1)$ to the normalized squared
value and derivative sums. In particular,
\begin{equation}\label{eq:non-tail-mass}
 S_x^{-2}\sum_{v\ \mathrm{non\text{-}tail}}|v(x)|^2=1-o(1).
\end{equation}
All these statements are uniform over $x\in\mathcal X_k$.

The value $U_x(0)$ is a standard real Gaussian. We verify that some
directional derivative has a uniformly positive conditional
variance given this value. Let $z$ be the coefficient-space vector
whose coordinates for packet $v$ are
$(\operatorname{Re}v(x),\operatorname{Im}v(x))$, and let $T$ be the
matrix whose corresponding columns are the four-dimensional real
derivative vectors $\dd_Y\operatorname{Re}v(0)$ and
$\dd_Y\operatorname{Im}v(0)$. Thus $|z|^2=S_x^2$. If $g$ is the
vector of independent standard Gaussian coefficients, the value and
gradient are $z^{\mathsf t}g/S_x$ and $Tg/S_x$. The conditional
gradient covariance is
\begin{equation}\label{eq:conditional-covariance}
 S_x^{-2}T\left(I-\frac{zz^{\mathsf t}}{S_x^2}\right)T^{\mathsf t}.
\end{equation}
For each non-tail packet define the unit coefficient vector $c_v$
supported on its two coordinates, with entries
\[
 \frac{(-\operatorname{Im}v(x),\operatorname{Re}v(x))}{|v(x)|}.
\]
These vectors are mutually orthonormal and perpendicular to $z$.
The orthogonal projection in \eqref{eq:conditional-covariance}
therefore dominates $\sum_v c_vc_v^{\mathsf t}$ as a nonnegative
quadratic form. Moreover
\[
 Tc_v=|v(x)|\operatorname{Im}\frac{\dd_Yv}{v}(0).
\]
Equations~\eqref{eq:local-imaginary-lower} and
\eqref{eq:non-tail-mass} show that the trace of
\eqref{eq:conditional-covariance} is at least a fixed $c>0$.
There is consequently a deterministic unit direction $e_x$ in
$\R^4$ such that the derivative
$Z_x=\partial_{e_x}U_x(0)$, conditioned on $U_x(0)$, has variance
between fixed positive constants. Its unconditional variance is
bounded above by the deterministic squared-derivative estimates.
Those estimates also bound the regression coefficient
$\operatorname{Cov}(Z_x,U_x(0))$.

For clarity, the needed uniform control of the local $C^2$ norm
follows from finitely many derivatives. The deterministic estimates
give
\[
 \sup_{|Y|\le2}\E|D_Y^qU_x(Y)|^2\le C_q
 \qquad(0\le q\le5).
\]
Integration and the Sobolev inequality on concentric balls in
$\R^4$ yield
\begin{equation}\label{eq:local-c2-moment}
 \E\lVert U_x\rVert_{C^2(\overline{B(0,1)})}^2
 \le C\E\lVert U_x\rVert_{H^5(B(0,2))}^2\le C.
\end{equation}
We choose the $C^2$ norm to control the Euclidean gradient norm and
Hessian operator norm, changing only the fixed constant in this
inequality. The Sobolev embedding used here is valid since
$5>2+4/2$; see \cite[Sections~5.6--5.7]{Evans2010}.

Let $a'=1/4$ and fix $0<c'<a'/4$. For every sufficiently large
fixed number $H$, the Gaussian conditional density formula gives
\[
 \Pp\bigl(1\le Z_x\le2,\ |U_x(0)|<c'/H\bigr)\ge c_1/H.
\]
Indeed, on the indicated value interval the conditional mean is
bounded, and the conditional variance is bounded above and away
from zero. The Gaussian density on the fixed derivative interval
$[1,2]$ then has a uniform positive lower bound. By
\eqref{eq:local-c2-moment}, failure of the bound
$\lVert U_x\rVert_{C^2}\le H$ has probability at most $C/H^2$.
Choose $H$ once so large that the intersection of these three
conditions has probability at least
$\pi_*=c_1/(2H)>0$.

On that intersection Taylor's formula at $\pm a'e_x/H$ gives
opposite signs, with absolute values at least $a'/(2H)$:
the linear term has magnitude at least $a'/H$, the center value
has magnitude at most $c'/H$, and the quadratic error is at most
$a'^2/(2H)$. The gradient bound preserves these signs on the
parallel transverse disks of radius
\[
 \sigma_* = a'/(4H^2)
\]
centered at those two points. The whole closed cylinder between
the disks is compactly contained in the unit ball. The direction
$e_x$, centers and radii are deterministic functions of the packet
data. After scaling by $R_x$, they provide a potential certificate
pair inside $B(x,R_x)$.

Let $F_k$ be the fraction of successful balls weighted by $R_x^3$.
No independence among these successes is claimed or needed.
Since $0\le F_k\le1$ and $\E F_k\ge\pi_*$,
\begin{equation}\label{eq:weighted-success}
 \Pp(F_k\ge\pi_*/2)\ge\frac{\pi_*}{2-\pi_*}>0.
\end{equation}
The Euclidean three-area of one successful disk is
$\omega_3\sigma_*^3R_x^3$, where $\omega_3$ is the unit-ball
volume in $\R^3$. On the event in \eqref{eq:weighted-success},
\eqref{eq:weighted-packing} therefore gives a certificate of size
at least $c_*k^{1+\gamma}$. Its cylinders lie in disjoint open
balls, as required by Definition~\ref{def:certificate}. This proves
the lemma with $p_* =\pi_* /(2-\pi_*)$.
\end{proof}

Lemmas~\ref{lem:noncancellation} and~\ref{lem:signs} can be realized
simultaneously. For each fixed choice of the orders, first fix the
positive number $p_*$ in Lemma~\ref{lem:signs}, and then take $k$
so large that the failure probability in
\eqref{eq:noncancellation} is less than $p_*/2$. The simultaneous
event has probability at least $p_*/2$. Thus, for every sufficiently
large allowed $k$, there is a real smooth $u$ satisfying
\eqref{eq:sum-weighted-bounds}, the pointwise lower bound in
\eqref{eq:noncancellation}, and a certificate of size
$c_*k^{1+\gamma}$. This probability subtraction is valid even if
$p_*$ is extremely small. The probability bounds above were obtained
without conditioning on $\mathcal B_k$; only deterministic derivative
estimates and interpolation used that event.

\subsection{Exact coefficient correction}

\begin{lemma}\label{lem:correction}
Let $(A,\rho)$ be a fixed smooth positive pair, let $u$ be real and
smooth, and let $W$ be a smooth positive weight. Suppose that
\[
 |u|+|\dd u|_{g_0}\ge k^{-500}W.
\]
Fix an integer $R\ge0$. Suppose that
\[
 |D^qu|\le C_qk^{2q+41}W\quad(0\le q\le R+2),\qquad
 |D^q\mathfrak r|\le C_qk^{-M}W\quad(0\le q\le R+1),
\]
where $\mathfrak r=\diver_\mu(A\,\dd u)+k^2\rho u$.
There are smooth real corrections $(\delta A,\delta\rho)$ canceling
the eigenfunction residual exactly, and
\begin{equation}\label{eq:correction-size}
 \lVert(\delta A,\delta\rho)\rVert_{C^R}
 \le C_R k^{-M+2129+1088R}.
\end{equation}
The corrections vanish on every open set where $\mathfrak r=0$.
In particular, choosing $M\ge2131+1088R$ makes them tend to zero
in $C^R$ as $k\to\infty$. For sufficiently small corrections the
output pair is positive and $u$ is its exact eigenfunction at $k^2$.
\end{lemma}

\begin{proof}
Write $I=g_0^{-1}$ and define globally
\begin{equation}\label{eq:correction-formula}
 \begin{split}
 \mathcal D&=u^2+I(\dd u,\dd u),\qquad F=-\mathfrak r/\mathcal D,\\
 \delta A&=FuI,\qquad
 \delta\rho=k^{-2}\bigl(Fu-\diver_\mu(FI\,\dd u)\bigr).
 \end{split}
\end{equation}
The lower bound on the first jet gives
$\mathcal D\ge\tfrac12 k^{-1000}W^2>0$, so these are smooth.
The product rule gives the exact identity
\begin{align*}
 \diver_\mu(\delta A\,\dd u)+k^2\delta\rho\,u
 &=u\diver_\mu(FI\,\dd u)+F I(\dd u,\dd u)\\
 &\quad+Fu^2-u\diver_\mu(FI\,\dd u)\\
 &=F\mathcal D=-\mathfrak r.
\end{align*}
Thus $\diver_\mu((A+\delta A)\dd u)
+k^2(\rho+\delta\rho)u=0$ identically.

We include finite derivative bounds to specify the accuracy needed
in this step. For $0\le q\le R+1$, the hypotheses imply
\[
 |D^q\mathcal D|\le C_q k^{2q+86}W^2.
\]
Every term in a derivative of order $q$ of $\mathcal D^{-1}$ is a
constant times
$\mathcal D^{-m-1}\prod_{i=1}^m D^{q_i}\mathcal D$, where
$q_i\ge1$, $\sum_iq_i=q$, and $m\le q$; the order-zero case is
included separately. Hence
\[
 |D^q(\mathcal D^{-1})|
 \le C_q k^{1000+2q+1086q}W^{-2}
 =C_q k^{1000+1088q}W^{-2}.
\]
By the product rule,
\begin{equation}\label{eq:F-derivatives}
 |D^qF|\le C_qk^{-M+1000+1088q}W^{-1}
 \qquad(0\le q\le R+1).
\end{equation}
No derivatives of $W^{-1}$ are required for these estimates: the
displayed weights bound each differentiated factor before the
factors are multiplied.

Multiplying \eqref{eq:F-derivatives} by derivatives of $u$ gives
$|D^q(Fu)|\le C_qk^{-M+1041+1088q}$. Replacing $u$ by $I\dd u$
gives the bound $C_qk^{-M+1043+1088q}$, with the derivatives of
the fixed tensor $I$ absorbed in the constant. The divergence in
\eqref{eq:correction-formula} costs one further derivative, and
the density correction has the factor $k^{-2}$. These bounds give
\eqref{eq:correction-size}; its exponent for the density is
$1043+1088(R+1)-2=2129+1088R$.
The cancellation of the two weights is what permits arbitrarily
small coefficient changes even where the original function is
exponentially small.

If $\mathfrak r=0$ on an open set, then $F=0$ there and both
corrections vanish. Finally, the cone of positive pairs is open
in $C^0$ on the compact base, so sufficiently small corrections
preserve positivity.
\end{proof}

\subsection{Simplicity on a protected round patch}

\begin{lemma}\label{lem:simplicity}
Suppose a smooth positive pair has a real eigenfunction $u$ at
$k^2=N(N+3)$, with $N>1$. Suppose there is an open region on which
the pair equals the round pair and $u=u_*$, containing a point with
$0<t<1$ at which $\dd u$
and $\dd t$ are independent. Then one can make $k^2$ simple by
an arbitrarily small smooth tensor perturbation supported in that
region, while preserving $u$ and $\rho$ exactly. The perturbation
can be arbitrarily small in any prescribed finite $C^R$ norm.
\end{lemma}

\begin{proof}
Take a small product coordinate patch $O$ in the stated region on
which $(u,t)$ are the first two coordinates, and $\dd u\ne0$.
Choose a smooth nonnegative bump $\chi$ supported in $O$ and
strictly positive on a smaller product patch $O_0$. Define
\[
 E=\chi\left(I-
 \frac{(I\dd u)\otimes(I\dd u)}{I(\dd u,\dd u)}\right),
 \qquad I=g_0^{-1},
\]
extending by zero off $O$. This is a smooth nonnegative symmetric
contravariant tensor and $E\dd u=0$.

Let $\mathcal E$ be the real eigenspace at $k^2$ for the given
pair. The nonnegative form
\[
 f\longmapsto\int_B E(\dd f,\dd f)\,\dd\mu,
 \qquad f\in\mathcal E,
\]
has kernel exactly $\operatorname{span}\{u\}$. Indeed, a kernel
element satisfies $\dd f\in\operatorname{span}\{\dd u\}$ on
$O_0$. After making the product patch smaller, this means
$f=b(u)$ for a smooth one-variable function $b$. The round
eigenfunction equations imply
\begin{equation}\label{eq:simplicity-composition}
 b''(u)|\dd u|_{g_0}^2=k^2\bigl(ub'(u)-b(u)\bigr).
\end{equation}
The Euclidean gradient of the homogeneous harmonic polynomial
defining $u_*$ has squared norm $N^2(1-t)^{N-1}$, and its radial
derivative is $Nu$. Tangential projection onto $S^4$ consequently
gives
\[
 |\dd u|_{g_0}^2=N^2\bigl((1-t)^{N-1}-u^2\bigr).
\]
Differentiate \eqref{eq:simplicity-composition} with respect to $t$
while fixing the coordinate $u$. Since $N>1$ and $t<1$, it follows
that $b''(u)=0$. Substitution back into
\eqref{eq:simplicity-composition} eliminates the constant term of
$b$, so $f=cu$ on $O_0$.

The difference $f-cu$ satisfies a second-order elliptic equation
with smooth real positive definite principal coefficients and
smooth lower-order coefficients. Aronszajn's unique continuation
theorem applies: its local hypotheses allow principal coefficients
of class $C^2$ with Lipschitz second derivatives and bounded
lower-order coefficients, and smooth solutions satisfy its
solution hypotheses \cite[pp.~235--236]{Aronszajn1957}. Vanishing
on an open patch gives infinite-order vanishing there. On the
connected base the difference is therefore identically zero,
proving the assertion about the kernel.

For $\tau>0$ replace $A$ by $A+\tau E$. The eigenfunction $u$ and
its eigenvalue remain exact because $E\dd u=0$ identically. We
claim that $k^2$ is simple for every sufficiently small positive
$\tau$. Otherwise there would be $\tau_m\downarrow0$ and real
eigenfunctions $f_m$ at $k^2$, normalized in $L^2(\rho\,\dd\mu)$
and orthogonal to $u$. Their energy identity is
\[
 \int_B(A+\tau_m E)(\dd f_m,\dd f_m)\,\dd\mu=k^2,
\]
so uniform ellipticity bounds them in $H^1$. A subsequence
converges weakly in $H^1$ and strongly in $L^2$ to a normalized
function $f_0$ orthogonal to $u$. Passing to the weak equation
shows $f_0\in\mathcal E$; elliptic regularity makes it smooth.
For any fixed $\phi\in\mathcal E$, testing the two eigenfunction
equations against one another cancels the original form and gives
\[
 \int_B E(\dd f_m,\dd\phi)\,\dd\mu=0.
\]
Weak convergence passes this fixed bilinear identity to the limit.
Now choose $\phi=f_0$ to see that $f_0$ is in the kernel just
identified, contradicting its normalization and orthogonality.
This proves the claim. The tensor $E$ is fixed before choosing
$\tau$; hence $\tau E$ can meet every prescribed finite norm
tolerance, while remaining supported in $O$.
\end{proof}

\subsection{Proof of the quantitative step}

\begin{proof}[Proof of Proposition~\ref{prop:step}]
Fix the input pair, the collar, $r$ and $\eta$. Choose
$R\ge\max\{r,1\}$ and a smaller positive tolerance $\eta'$ so
that total $C^R$ movement less than $\eta'$ implies the requested
$C^r$ movement and keeps the pair positive and in $\mathcal U$.
Such a choice exists because the fixed input belongs to the open
set $\mathcal U$.

Choose $J=R+2$ and $M\ge2131+1088R$, and then choose the finite
packet orders giving \eqref{eq:sum-weighted-bounds}. All these
choices precede frequency selection. Lemmas~\ref{lem:noncancellation}
and~\ref{lem:signs} give, for every sufficiently large allowed $k$,
a sample with the required first-jet bound, residual bounds and a
certificate of size $c_*k^{1+\gamma}$. Lemma~\ref{lem:correction}
changes the pair by at most $C_Rk^{-2}$ in $C^R$ while preserving
that same function. Choose $k$ large enough for this change to be
less than $\eta'/2$.

We describe explicitly why a narrow collar only changes this
frequency threshold. Packet supports avoid the region below
$s_1-Ck^{-1/3-100\delta}$. Requiring
\begin{equation}\label{eq:collar-thresholds}
 Ck^{-1/3-100\delta}<w/100,
 \qquad
 cw k>(M'+P)\log k+\log C_q(w)
\end{equation}
protects a fixed subcollar and makes any required cutoff residual
$C_q(w)k^Pe^{-cw k}$ smaller than $k^{-M'}$. Only finitely many
orders $q$ and finite targets $M'$ are needed at this step.
For every already fixed $w>0$, all inequalities in
\eqref{eq:collar-thresholds} hold for every sufficiently large $k$.
In particular the exact corrections vanish below $s_1-w/100$,
where the coefficients remain round and $u=u_*$.

In the subcollar
\[
 l_-+w/20<t<l_-+w/10
\]
choose a point with $\sin(N\vartheta)\ne0$. Since $0<t<1$,
$\dd t\ne0$, and the nonzero angular derivative makes $\dd u$
and $\dd t$ independent. Lemma~\ref{lem:simplicity} applies there.
After $k$ and the sample are fixed, form its tensor $E$ and choose
$\tau>0$ below its simplicity threshold and so small that
$\lVert\tau E\rVert_{C^R}<\eta'/2$. This tensor can have large
finite derivatives depending on $k$; the subsequent scalar choice
of $\tau$ absorbs them. It preserves $u$, its exact eigenvalue and
all current signs exactly. Its support lies in $t>l_-$, so the
final pair is round on $t<l_-$.

The final pair meets every conclusion of
Proposition~\ref{prop:step}. In particular, its certificate constant
$c_*$ was fixed after choosing the finite orders and before taking
$k$ large, and it was unaffected by either correction. The
exponent $\gamma$, the chart and the background neighborhood were
fixed upstream. Neither the collar estimates nor the arbitrary
finite tolerances change them. The construction works for every
sufficiently large integer $N$, as asserted.
\end{proof}

\subsection{Spectral persistence under finite norm tolerances}

\begin{lemma}\label{lem:persistence}
Let a smooth positive pair $\mathcal P=(A,\rho)$ have a simple
positive base eigenvalue $\lambda$ and a real eigenfunction with a
sign certificate $\mathcal C$. If $I\subset(0,\infty)$ is an open
interval containing $\lambda$, there are an integer $s\ge1$ and
$\varepsilon>0$ such that every smooth positive pair
$\mathcal P'$ satisfying
\[
 \lVert\mathcal P'-\mathcal P\rVert_{C^s}\le\varepsilon
\]
has a simple eigenvalue in $I$ and a real eigenfunction with the
same certificate disks. Thus a closed ball defined by one finite
norm can be contained in a persistence neighborhood.
\end{lemma}

\begin{proof}
We recall the spectral and regularity facts used here. The form
\[
 \int_B A(\dd f,\dd v)\,\dd\mu+
 \int_B\rho fv\,\dd\mu
\]
is coercive on $H^1(B)$ and defines the inverse of $-L+1$ in
$L^2(\rho\,\dd\mu)$. The compact embedding
$H^1\hookrightarrow L^2$ makes this inverse compact, and it is
selfadjoint. The spectral theorem and elliptic regularity give a
discrete nonnegative spectrum with finite multiplicities and smooth
real eigenfunctions. Its ordered eigenvalues obey the min--max
principle with Rayleigh quotient
\begin{equation}\label{eq:weighted-rayleigh}
 \frac{\int_B A(\dd f,\dd f)\,\dd\mu}
      {\int_B\rho f^2\,\dd\mu}.
\end{equation}

Uniform multiplicative comparisons of nearby positive tensors and
densities compare \eqref{eq:weighted-rayleigh} for every $f$ and
therefore imply convergence of each ordered eigenvalue under
$C^0$ convergence of pairs. Choose a smaller interval about
$\lambda$, with closure in $I$, separating it from its neighboring
ordered eigenvalues. Their convergence implies that for sufficiently
close pairs the corresponding ordered eigenvalue is still simple
and lies in $I$.

We also need uniform convergence of a corresponding normalized
eigenfunction, up to sign. Here only finitely many coefficient
derivatives are needed. For a fixed smooth elliptic $L$ on the
closed base, the local interior estimates, a finite partition of
unity and interpolation give
\begin{equation}\label{eq:elliptic-persistence-estimate}
 \lVert f\rVert_{H^{a+2}}
 \le C_a\bigl(\lVert Lf\rVert_{H^a}+\lVert f\rVert_{L^2}\bigr)
 \qquad(a=0,1,2,\ldots).
\end{equation}
These are the usual interior elliptic estimates; the Sobolev and
compactness statements used here are described in
\cite[Sections~5.6--5.7 and~6.3]{Evans2010}. A sufficiently small
finite coefficient norm makes $L'-L$ small as an operator
$H^{a+2}\to H^a$ for each of finitely many specified $a$. Its
contribution in \eqref{eq:elliptic-persistence-estimate} can then
be absorbed on the left. Applying the resulting estimates at
$a=0,2$, together with the eigenfunction equation, gives a uniform
$H^4$ bound for eigenfunctions whose eigenvalues stay in the
chosen bounded interval and whose $L^2$ norms are bounded.
For example, requiring sufficiently small $C^4$ changes of the
pair is more than enough for these finitely many operator bounds:
the principal coefficients of $L$ involve $A,\rho$ without
derivatives, and its first-order coefficients involve at most
their first derivatives.

Suppose that pairs converge in such a finite norm to $\mathcal P$.
Normalize their corresponding real eigenfunctions $f_m$ by
$\int\rho_m f_m^2\,\dd\mu=1$. Positivity bounds their ordinary
$L^2$ norms, so the preceding argument bounds them in $H^4$.
Rellich compactness and Sobolev embedding in dimension four give
a subsequence converging in $C^0$ (for example, first converge in
$H^3$, which embeds in $C^0$). Passing to the weak equation shows
that the limit is a normalized eigenfunction at $\lambda$ for
$\mathcal P$. By simplicity it equals the prescribed normalized
eigenfunction or its negative.

The finite collection of strict signs on compact disks has a
positive minimum absolute margin after this normalization. Hence
all sufficiently close such eigenfunctions have the same opposite
signs on the same disks, allowing an overall change of sign. If
no sufficiently small finite-norm neighborhood had this property,
a sequence converging in that norm would contradict the preceding
compactness argument. Finally, decrease its radius to put a closed
finite-norm ball inside the neighborhood. This proves the lemma,
including the assertion for varying densities.
\end{proof}

\subsection{One smooth limiting pair and the ordinary metric}

\begin{proof}[Proof of Theorem~\ref{thm:main}]
Use monotone fixed coordinate norms for pairs. Choose $b>0$ so
that the closed $C^1$ ball of radius $b$ about
$\mathcal P_0=(g_0^{-1},1)$ consists of positive pairs in
$\mathcal U$. Fix
\[
 l^{(j)}=\frac14+2^{-j-2}\quad(j\ge0),\qquad
 \epsilon_0=\gamma/4>0.
\]
We construct pairs $\mathcal P_j=(A_j,\rho_j)$, frequencies $k_j$,
and certificates $\mathcal C_j$ inductively. Every output is round
on $t<l^{(j)}$.

After producing $\mathcal P_i$, apply Lemma~\ref{lem:persistence}
to its simple eigenvalue $k_i^2$, its certificate, and the interval
$(k_i^2/2,2k_i^2)$. Denote a resulting finite norm order and
closed-ball radius by $s_i$ and $\varepsilon_i>0$.
At stage $j$, these numbers are already fixed for $i<j$. Put
\begin{equation}\label{eq:diagonal-budgets}
 \begin{split}
 r_j&=\max\{j+1,1,s_1,\ldots,s_{j-1}\},\\
 \eta_j&=\min\left\{2^{-j},\ b2^{-j-1},\,
 \min_{1\le i<j}2^{-(j-i+1)}\varepsilon_i\right\}.
 \end{split}
\end{equation}
For $j=1$, the inner minimum is omitted. These are finite
requirements at each stage. Apply Proposition~\ref{prop:step}
to $\mathcal P_{j-1}$, with
$l_+=l^{(j-1)}$, $l_-=l^{(j)}$, order $r_j$ and tolerance
$\eta_j$. Use the certificates from its construction in
Lemma~\ref{lem:signs}, with balls centered in $V$ and radii at most
$k_j^{-1}$. The proposition gives a positive constant $c_j$ and all
sufficiently large allowed frequencies. Choose one with
$k_j>2k_{j-1}$ when $j>1$, and also so large that
\[
 c_jk_j^{\gamma/2}>j.
\]
The output certificate then satisfies
\begin{equation}\label{eq:stage-certificate-size}
 |\mathcal C_j|>j k_j^{1+\gamma/2}.
\end{equation}
Only after this output has been constructed do we select its own
$s_j,\varepsilon_j$ by Lemma~\ref{lem:persistence}. They constrain
subsequent stages, so no current frequency choice depends
circularly on its own future persistence radius.

The increments are summable in every fixed $C^r$ norm: for all
large $j$, $r_j\ge r$ and their $C^{r_j}$ norms are less than
$2^{-j}$. Therefore the pairs converge to one smooth pair
$\mathcal P_\infty=(A_\infty,\rho_\infty)$. Their total $C^1$
movement from the round pair is less than
$b\sum_{j\ge1}2^{-j-1}=b/2$, so the limit remains positive and
belongs to $\mathcal U$. For each fixed $i$, the same budgets give
\[
 \lVert\mathcal P_\infty-\mathcal P_i\rVert_{C^{s_i}}
 \le\sum_{j>i}2^{-(j-i+1)}\varepsilon_i
 =\varepsilon_i/2.
\]
The limit lies in every selected persistence neighborhood. Hence
there are real smooth eigenfunctions $v_i$ for the limiting base
operator, with the same certificate disks $\mathcal C_i$, and
\begin{equation}\label{eq:limiting-eigenvalue-intervals}
 -L_\infty v_i=\lambda_i v_i,\qquad
 k_i^2/2<\lambda_i<2k_i^2.
\end{equation}
The frequency separation makes these eigenvalue intervals disjoint
and tending to infinity. The functions $v_i$ are nonzero by their
normalization, or just by their strict signs.

Set $G_\infty^{-1}=A_\infty/\rho_\infty$ and let
\[
 q=\frac{\rho_\infty\,\dd\mu}{\dd\operatorname{vol}_{G_\infty}}.
\]
This is a smooth positive function, the ratio of two smooth
positive base densities. We use the standard warped-product
realization of a weighted Laplacian
\cite[Section~4, equation~(2)]{CharalambousLu2015}, with a circle
as the fiber. On the fixed manifold $M=B\times S^1$
define
\begin{equation}\label{eq:circle-metric}
 g=G_\infty+q^2\dd\varphi^2,
\end{equation}
where $\dd\varphi^2$ is the standard unit-circle metric. The
manifold is smooth, closed and connected, of dimension five, and
\[
 \dd\operatorname{vol}_g
 =q\,\dd\operatorname{vol}_{G_\infty}\,\dd\varphi
 =\rho_\infty\,\dd\mu\,\dd\varphi.
\]
For a function $f$ pulled back from $B$, the coordinate divergence
formula gives
\begin{equation}\label{eq:circle-operator}
 \Delta_g f
 =\frac{1}{\rho_\infty\mu}
   \partial_a\bigl(\rho_\infty\mu
                    G_\infty^{ab}\partial_b f\bigr)
 =\rho_\infty^{-1}\diver_\mu(A_\infty\dd f)
 =L_\infty f.
\end{equation}
Here $\mu$ in the coordinate expression denotes its positive
coordinate density. Thus the pullbacks $u_i$ of $v_i$ are exact
real Laplace--Beltrami eigenfunctions with eigenvalues
$\lambda_i$.

It remains to convert their certificates into the claimed measure
bound. Fix a closed interval $I_0$ of positive length in a circle
coordinate chart. All certificate cylinders are in one fixed
compact subset of $\Omega$: their centers lie in $V\Subset\Omega$
and their radii tend uniformly to zero, with the initial
frequencies chosen large enough for the same inclusion. On the
product of this compact set with $I_0$, coordinate and Riemannian
distances have fixed Lipschitz comparison constants. One may see
this directly by extending the coordinate functions smoothly with
cutoffs to the compact manifold; their gradients are then
uniformly bounded in the single fixed metric $g$.

For one certificate cylinder $C_i^{(a)}$, let $D_i^{(a)}$ be one
of its transverse disks. The opposite endpoint signs imply that
every longitudinal segment through the disk contains a zero of
$v_i$. Therefore the subset
\[
 \Z(u_i)\cap(C_i^{(a)}\times I_0)
\]
projects onto $D_i^{(a)}\times I_0$ by coordinate orthogonal
projection in the base and the circle coordinate. The Lipschitz
constant of every such projection is bounded by one fixed $K$,
independent of the disk's radius and orientation. The defining
covering inequality for Hausdorff measure under a Lipschitz map
gives
\[
 \Hh_g^4\bigl(\Z(u_i)\cap(C_i^{(a)}\times I_0)\bigr)
 \ge K^{-4}\Hh_{\mathrm{Euc}}^4(D_i^{(a)}\times I_0)
 =K^{-4}|I_0|\Hh_{\mathrm{Euc}}^3(D_i^{(a)}).
\]
The closed cylinders lie in disjoint open balls. Their indicated
nodal subsets are consequently pairwise disjoint Borel sets, so
their Hausdorff measures add. Using
\eqref{eq:stage-certificate-size}, we obtain one constant $c>0$
belonging to the final metric such that
\begin{equation}\label{eq:final-nodal-lower}
 \Hh_g^4(\Z(u_i))\ge c|\mathcal C_i|
 > c i k_i^{1+\gamma/2}.
\end{equation}
This argument uses continuity and signs only, without any
regularity assumption on the nodal set.

Finally, \eqref{eq:limiting-eigenvalue-intervals} and
$\epsilon_0=\gamma/4$ imply
\[
 \lambda_i^{1/2+\epsilon_0}
 \le2^{1/2+\epsilon_0}k_i^{1+\gamma/2}.
\]
Combining this with \eqref{eq:final-nodal-lower} yields
\[
 \frac{\Hh_g^4(\Z(u_i))}{\lambda_i^{1/2+\epsilon_0}}
 \ge c\,2^{-1/2-\epsilon_0}i\longrightarrow\infty.
\]
Both $g$ and $\epsilon_0$ are fixed throughout this sequence,
which proves Theorem~\ref{thm:main}.
\end{proof}

\begin{corollary}
For every integer $n>5$ there is a smooth closed connected
$n$-dimensional Riemannian manifold and a fixed $\epsilon_0>0$
with the same divergence of nodal measure divided by
$\lambda^{1/2+\epsilon_0}$.
\end{corollary}

\begin{proof}
Take the product of \eqref{eq:circle-metric} with any fixed smooth
closed connected manifold of dimension $n-5$, and pull back the
same eigenfunctions. The product Laplacian preserves their
eigenvalues. Multiply each of the projection certificates above
by a fixed compact coordinate box of positive $(n-5)$-volume in
the extra factor. The same Lipschitz projection argument, now in
dimension $n-1$, gives a fixed positive multiple of the preceding
certificate lower bound. The denominator and its fixed exponent
are unchanged, so the ratio still tends to infinity.
\end{proof}
